\chapter*{\centering \large Resumen}
\addcontentsline{toc}{chapter}{Resumen}
\markboth{Resumen}{Resumen}

Esta tesis estudia la convergencia de aproximaciones de la ecuación de Poisson mediante redes neuronales físicamente informadas (PINNs), con énfasis en la relación entre la pérdida residual y el error de solución. Se reconstruye el marco de Shin, Darbon y Karniadakis para problemas elípticos, precisando las hipótesis de regularidad, muestreo y estabilidad, y se estudia una extensión condicional con comparadores y minimizadores aproximados, cuyas tolerancias deben estar controladas por la escala del peso regularizador. El análisis distingue las garantías referidas a minimizadores globales de los resultados obtenidos mediante optimización numérica finita. Se estudian problemas unidimensionales, se examina la reproducción de experimentos de referencia y se resuelve Poisson en dos dimensiones con condiciones de Dirichlet impuestas mediante penalización o mediante una parametrización que las satisface exactamente. Para el problema bidimensional se realizan veinticuatro entrenamientos, correspondientes a cuatro resoluciones, tres semillas y dos tratamientos de frontera. En la resolución más alta, la mediana del error relativo $L^2$ disminuye de $6.19\times10^{-3}$ a $1.74\times10^{-4}$ al utilizar Dirichlet fuerte. La evaluación independiente de residuos y errores muestra, además, que una pérdida reducida en los puntos de colocación puede coexistir con residuos mayores fuera de ellos. La reproducción del caso oscilatorio presenta discrepancias que no se atribuyen a una causa única. Un estudio adicional con redes $\tanh$ de dos capas ocultas distingue continuación de la optimización (A), anchura con muestreo fijo (B) y escalamiento conjunto de muestras y anchura (C). La pérdida residual disminuye, pero en C la pendiente log--log descriptiva ($-0.284$, tres escalas) es menos pronunciada que la escala suficiente de referencia ($-0.750$); el presupuesto de optimización no es proporcional al tamaño de muestra. Los resultados respaldan la utilidad de imponer exactamente las condiciones de frontera en el problema estudiado, pero no constituyen una verificación de las hipótesis asintóticas ni una comparación general de superioridad frente a los métodos de elementos y volúmenes finitos.

\noindent\textbf{Palabras clave:} redes neuronales físicamente informadas, ecuación de Poisson, convergencia, condiciones de Dirichlet, error residual.

\clearpage
\chapter*{\centering \large Abstract}
\addcontentsline{toc}{chapter}{Abstract}
\markboth{Abstract}{Abstract}

This thesis studies the convergence of approximations to the Poisson equation using physics-informed neural networks (PINNs), with particular attention to the relationship between residual loss and solution error. The convergence framework of Shin, Darbon, and Karniadakis for elliptic problems is reconstructed, making explicit the regularity, sampling, and stability assumptions required to relate an empirical loss to functional error estimates. A conditional extension admitting approximate comparators and approximately optimal solutions is examined, subject to quantitative tolerances on the scale of the regularization weight. The analysis distinguishes theoretical guarantees for global minimizers from results produced by finite numerical optimization. One-dimensional problems are examined, reference experiments are revisited, and the two-dimensional Poisson equation is solved with Dirichlet conditions imposed either by a penalty or by a parameterization that satisfies them exactly. For the two-dimensional problem, twenty-four training runs are performed using four resolutions, three random seeds, and two boundary treatments. At the finest resolution, the median relative $L^2$ error decreases from $6\mathord{.}19\times10^{-3}$ to $1\mathord{.}74\times10^{-4}$ when strong Dirichlet enforcement is used. Independent evaluations of residuals and errors also reveal that small losses at collocation points may coexist with substantially larger residuals away from those points. The oscillatory reference example exhibits discrepancies that cannot be assigned to a single established cause. Additional experiments with two-hidden-layer $\tanh$ networks distinguish optimization continuation (A), width variation with fixed samples (B), and joint sample--width scaling (C). Although residual losses decrease, the descriptive log--log slope in C ($-0.284$, three scales) is shallower than the sufficient reference exponent ($-0.750$); the optimization budget does not scale proportionally with sample count. These findings support the usefulness of exactly imposing boundary conditions in the problem studied, but they do not verify the asymptotic assumptions of the convergence theorems or establish general superiority over finite element or finite volume methods.

\noindent\textbf{Keywords:} physics-informed neural networks, Poisson equation, convergence, Dirichlet boundary conditions, residual error.